\DocumentMetadata{
  pdfversion=2.0,
  pdfstandard=ua-2,
  lang=en-US,
  tagging=on,
  tagging-setup={math/setup=mathml-SE,tabsorder=structure}
}
\documentclass[11pt,letterpaper]{article}
\usepackage[margin=0.68in,includefoot]{geometry}
\usepackage{newcomputermodern}
\usepackage{amsmath,amscd,mathtools}
\usepackage{tikz,adjustbox}
\providecommand{\square}{\mdlgwhtsquare}
\usepackage[hidelinks]{hyperref}
\hypersetup{
  pdftitle={Algebra First-Year Exam, January 2013, Part 1},
  pdfauthor={University of Florida Department of Mathematics},
  pdfsubject={Graduate mathematics examination},
  pdfdisplaydoctitle=true,
  bookmarksopen=true
}
\setcounter{secnumdepth}{0}
\setlength{\parindent}{0pt}
\setlength{\parskip}{0.28em}
\setlength{\emergencystretch}{3em}
\setlength{\leftmargini}{2.15em}
\setlength{\leftmarginii}{2.65em}
\setlength{\leftmarginiii}{3em}
\pagestyle{empty}
\raggedbottom
\allowdisplaybreaks
\NewTaggingSocketPlug{block/list/label}{examproblem}{%
  \tagstructbegin{tag=Lbl}\tagstructend%
  \tagstructbegin{tag=\LBody}%
  \tagstructbegin{tag=\ProblemHeadingTag,title-o={\ProblemHeadingText}}%
  \tagmcbegin{tag=\ProblemHeadingTag,actualtext={\ProblemHeadingText}}%
  #2%
  \tagmcend%
  \tagstructend%
}
\newenvironment{examproblems}[2]{%
  \def\ProblemHeadingTag{#2}%
  \AssignTaggingSocketPlug{block/list/label}{examproblem}%
  \begin{enumerate}%
  \setcounter{enumi}{#1}%
  \renewcommand{\labelenumi}{\textbf{\theenumi.}}%
  \setlength{\topsep}{0.35em}%
  \setlength{\partopsep}{0pt}%
  \setlength{\itemsep}{0.48em}%
  \setlength{\parsep}{0.18em}%
}{\end{enumerate}}
\newenvironment{examsubparts}{%
  \AssignTaggingSocketPlug{block/list/label}{default}%
  \begin{enumerate}%
  \setlength{\topsep}{0.18em}%
  \setlength{\partopsep}{0pt}%
  \setlength{\itemsep}{0.16em}%
  \setlength{\parsep}{0.1em}%
}{\end{enumerate}}
\newenvironment{examinstructions}{%
  \par\begingroup\itshape
}{%
  \par\endgroup\vspace{0.18em}
}
\makeatletter
\renewcommand\subsection{\@startsection{subsection}{2}{0pt}%
  {-1.25ex plus -.25ex minus -.1ex}%
  {0.55ex plus .1ex}%
  {\normalfont\large\bfseries}}
\renewcommand{\@maketitle}{%
  \newpage\null\vskip -1em%
  {\footnotesize\bfseries
    \href{https://www.ufl.edu/}{University of Florida}, \href{https://math.ufl.edu/}{Department of Mathematics}\par}%
  \vskip -0.5em%
  \tagstructbegin{tag=H1,title={Algebra First-Year Exam, January 2013, Part 1}}%
  \tagpdfparaOff%
  \tagmcbegin{tag=H1}%
    {\LARGE\bfseries\href{https://gma.math.ufl.edu/exams/algebra-first-year/}{Algebra First-Year Exam}, January 2013, Part 1\par}%
  \tagmcend%
  \tagpdfparaOn%
  \tagstructend%
  \vskip 0.25em%
  \tagmcbegin{artifact}%
    \hrule height 1pt%
  \tagmcend%
  \vskip 1em%
}
\makeatother
\title{Algebra First-Year Exam, January 2013, Part 1}
\author{}
\date{}
\begin{document}
\maketitle
\thispagestyle{empty}
\begin{examinstructions}
Answer four problems. Write your answers clearly in complete English sentences. You may quote results (within reason) as long as you state them clearly.
\end{examinstructions}
\begin{examproblems}{0}{H2}
\def\ProblemHeadingText{Problem 1.}
\item
Give an example of each of the following, with some explanation.
\begin{examsubparts}
\item[\textnormal{(a)}]
A group~\(G\) and a subgroup \(H \leq G\) such that \(C_G(H) \neq N_G(H)\).
\item[\textnormal{(b)}]
A group~\(G\) and a subgroup \(H \leq G\) such that \(H \nsubseteq C_G(H)\).
\item[\textnormal{(c)}]
A group~\(G\) and a normal subgroup \(H \trianglelefteq G\) such that \(C_G(H) \neq G\).
\end{examsubparts}
\def\ProblemHeadingText{Problem 2.}
\item
Let~\(G\) be a group and let \(T_G = \{x \in G : x^n = 1 \text{ for some } n \geq 1\}\).
\begin{examsubparts}
\item[\textnormal{(a)}]
Prove that if~\(G\) is abelian then \(T_G\) is a subgroup of~\(G\).
\item[\textnormal{(b)}]
Give an example of a nonabelian group~\(G\) such that \(T_G\) is not a subgroup of~\(G\). Justify any claims that you make.
\end{examsubparts}
\def\ProblemHeadingText{Problem 3.}
\item
Let~\(G\) be a finite group which acts on the set~\(S\). Let~\(H\) be a normal subgroup of~\(G\) with the following property: For every \(x_1, x_2 \in S\), there is a unique \(h \in H\) such that \(h(x_1) = x_2\). Choose \(x \in S\) and let \(G_x = \{g \in G : g(x) = x\}\) denote the stabilizer of~\(x\).
\begin{examsubparts}
\item[\textnormal{(a)}]
Prove that \(G = G_xH\) and \(H \cap G_x = 1\).
\item[\textnormal{(b)}]
Show that if \(H \leq Z(G)\) then \(G_x \trianglelefteq G\) and~\(G\) is the internal direct product of \(G_x\) and~\(H\).
\end{examsubparts}
\def\ProblemHeadingText{Problem 4.}
\item
This problem has two parts.
\begin{examsubparts}
\item[\textnormal{(a)}]
Prove that every simple subgroup of \(S_4\) is abelian.
\item[\textnormal{(b)}]
Use the result from (a) to prove that if~\(G\) is a nonabelian simple group then every proper subgroup of~\(G\) has index at least 5.
\end{examsubparts}
\def\ProblemHeadingText{Problem 5.}
\item
Let~\(G\) be a group of order \(76 = 4 \cdot 19\).
\begin{examsubparts}
\item[\textnormal{(a)}]
Prove that~\(G\) contains a normal Sylow 19-subgroup.
\item[\textnormal{(b)}]
Prove that the center of~\(G\) contains an element of order 2.
\end{examsubparts}
\def\ProblemHeadingText{Problem 6.}
\item
Let~\(G\) be a group.
\begin{examsubparts}
\item[\textnormal{(a)}]
Define the derived series \(G^{(0)} \geq G^{(1)} \geq G^{(2)} \geq \cdots\) of~\(G\).
\item[\textnormal{(b)}]
Prove that~\(G\) is solvable if and only if \(G^{(n)} = \{1\}\) for some \(n \geq 0\).
\item[\textnormal{(c)}]
Give an example of a solvable group~\(G\) such that \(G^{(1)} \neq \{1\}\).
\end{examsubparts}
\end{examproblems}
\end{document}
